% Main torque-control loop with short secondary model and null-space inputs.
\begin{tikzpicture}[
    font=\normalsize,
    blk/.style={draw=black, rounded corners=2pt,
                minimum height=10mm, inner xsep=2.5mm, align=center},
    src/.style={draw=black, rounded corners=2pt,
                minimum height=10mm, inner xsep=2.5mm, align=center},
    sel/.style={draw=black, rounded corners=2pt,
                minimum height=10mm, inner xsep=2.5mm, align=center},
    sig/.style={-{Latex[length=2mm]}, thick, black},
    lbl/.style={font=\small, inner sep=1.6pt},
  ]

% Main calculation path, read from left to right.
\node[src, text width=3.0cm] (model) at (0.0,0.0)
  {Robot Model\\[-2pt]\small
   $J(q),\ \tau_c(q,\dot q)$};
\node[blk] (error) at (4.5,0.0)
  {Cartesian Error\\[-2pt]\small $e_p,\ e_R$};
\node[blk, text width=4.6cm] (wrench) at (9.4,0.0)
  {\footnotesize Cartesian Impedance Wrench\\[-2pt]\small $F$};
\node[blk] (cart) at (14.3,0.0)
  {Cartesian Torque\\[-2pt]\small
   $\tau_{\mathrm{cart}}=J^{\!\top}F$};
\node[blk, text width=3.3cm] (sum) at (19.2,0.0)
  {Torque Sum\\[-2pt]\footnotesize
   $\tau_{\mathrm{cmd}}=\tau_{\mathrm{cart}}+\tau_{\mathrm{null}}+\tau_c$};

% Reference and secondary torque inputs.
\node[sel] (machine) at (9.4,2.25) {Controller State};
\node[sel] (null) at (19.2,2.25)
  {Null-Space Term\\[-2pt]\small $\tau_{\mathrm{null}}$};

% Physical actuation and measured joint feedback.
\node[blk] (motors) at (19.2,-4.25) {Joint Motors};
\node[src, text width=3.2cm] (feedback) at (0.0,-4.25)
  {Joint Feedback\\[-2pt]\small
   $q,\ \dot q,\ T_{EE}$\\[-2pt]
   $\dot p,\ \omega$};

% Main calculation path.  The labels are lifted clear of the short gaps while
% the signal arrows remain direct and terminate on the box boundaries.
\draw[sig] (error) -- node[lbl, above, yshift=5mm] {$e_p,\ e_R$} (wrench);
\draw[sig] (wrench) -- node[lbl, above, yshift=5mm] {$F$} (cart);
\draw[sig] (cart) -- node[lbl, above, yshift=5mm] {$\tau_{\mathrm{cart}}$} (sum);
\draw[sig] (sum.south east) --
      node[lbl, right] {$\tau_{\mathrm{cmd}}$} (sum.south east |- 0,-3.15) --
      (motors.north |- 0,-3.15) -- (motors.north);

% State-selected references and gains.
\draw[sig] (machine.west) -- node[lbl, above] {$p_d,\ R_d$}
      (4.5,2.25) -- (error.north);
\draw[sig] (machine.south) -- node[lbl, right, pos=0.45]
      {$K,\ D,\ \dot p_d$} (wrench.north);

% Model-dependent and secondary torque inputs.
\draw[sig] (model.north) -- (0.0,3.45) --
      node[lbl, above, pos=0.72] {$J(q)$} (14.3,3.45) -- (cart.north);
\draw[sig] (model.south east) -- (model.south east |- 0,-1.35) --
      node[lbl, above, pos=0.68] {$\tau_c(q,\dot q)$} (19.2,-1.35) --
      (sum.south);
\draw[sig] (null.south) -- node[lbl, right] {$\tau_{\mathrm{null}}$}
      (sum.north);

% Joint sensing closes the physical loop.
\draw[sig] (motors.west) -- node[lbl, above]
      {Joint motion} (feedback.east);
\draw[sig] (feedback.north) -- node[lbl, left]
      {Measured $q,\ \dot q$} (model.south);
\draw[sig] (feedback.north east) -- (feedback.north east |- 0,-3.35) --
      node[lbl, above] {$T_{EE}$} (2.95,-3.35) -- (2.95,0.28) --
      ($(error.west)+(0,0.28)$);
\draw[sig] (feedback.south) -- (0.0,-5.35) --
      node[lbl, below, pos=0.62] {$\dot p,\ \omega$} (9.4,-5.35) --
      (wrench.south);

\end{tikzpicture}
